\documentclass[11pt]{article}

\usepackage[a4paper,margin=1in]{geometry}
\usepackage[T1]{fontenc}
\usepackage[utf8]{inputenc}
\usepackage{lmodern}
\usepackage{amsmath,amssymb}
\usepackage[hidelinks]{hyperref}
\usepackage{enumitem}
\usepackage{xcolor}

\setlist[itemize]{topsep=4pt,itemsep=3pt}
\setlist[enumerate]{topsep=4pt,itemsep=3pt}

\title{arXiv Strategy for the Cut-and-Project Paper}
\author{Dirk Kunert}
\date{\today}

\begin{document}

\maketitle

\section*{Purpose}

This note collects practical suggestions for submitting the manuscript

\begin{center}
\emph{A Period Formula for Rational Cut-and-Project Strip Projections
with Multiplicity}
\end{center}

to arXiv.  It lists suitable subject categories, possible expert readers,
and a suggested first-contact email.

\section*{Can an independent researcher submit to arXiv?}

Yes.  An independent researcher may submit to arXiv.  The issue is not the
lack of a university affiliation, but the endorsement requirement.  A new
submitter normally needs endorsement for the relevant arXiv category or
subject class.

The affiliation may simply be written as

\begin{verbatim}
Independent Researcher
\end{verbatim}

in the paper.

\section*{Recommended arXiv categories}

The paper is not best framed primarily as deep number theory.  It uses
number-theoretic tools such as congruences, residue classes, and coprimality,
but its main object is a gap sequence arising from a geometric and
combinatorial construction.

Recommended categories:

\begin{itemize}
    \item Primary: \texttt{math.CO} or \texttt{math.DS}.
    \item Possible cross-listing: \texttt{math.NT} or \texttt{math.MG}.
\end{itemize}

A reasonable first choice is

\begin{center}
\texttt{math.CO}
\end{center}

because the proof is essentially combinatorial and residue-class based.
If the paper is presented more strongly as a symbolic-dynamical gap-sequence
result, then \texttt{math.DS} is also plausible.

\section*{How to approach potential endorsers or expert readers}

It is better not to send a cold message that immediately asks for endorsement.
A more professional strategy is:

\begin{enumerate}
    \item Send a concise email with the manuscript attached or linked.
    \item Explain the result in two or three sentences.
    \item Ask whether the result appears suitable as an arXiv preprint and
    whether the recipient knows prior literature containing the same formula.
    \item Only if the response is positive, ask whether the person would be
    willing to endorse the arXiv submission.
\end{enumerate}

This is less intrusive and more likely to receive a serious response.

\section*{Most relevant potential readers}

The following people are mathematically close to the topic.  They should not
be viewed as guaranteed endorsers; they are possible expert readers whose
published work is near cut-and-project sets, gap problems, symbolic dynamics,
or combinatorics on words.

\subsection*{Closest to cut-and-project sets and gap problems}

\begin{itemize}
    \item \textbf{Alan Haynes.}  Relevant because of work on gap problems and
    frequencies of patches in cut-and-project sets.

    \item \textbf{Henna Koivusalo.}  Relevant because of work on
    cut-and-project sets, gap problems, and aperiodic order.

    \item \textbf{Lorenzo Sadun.}  Relevant because of work on tilings,
    aperiodic order, and cut-and-project structures.

    \item \textbf{James Walton.}  Relevant because of work on aperiodic order,
    topology of tilings, and cut-and-project structures.

    \item \textbf{Reem Yassawi.}  Relevant because of work connecting symbolic
    dynamics, Meyer sets, and self-similar structures.

    \item \textbf{Val\'{e}rie Berth\'{e}.}  Relevant because of work on
    symbolic dynamics, combinatorics on words, numeration, and Meyer-set
    related structures.
\end{itemize}

\subsection*{Gap problems, lattices, and the three-gap theorem}

\begin{itemize}
    \item \textbf{Jens Marklof.}  Relevant because of work on the three-gap
    theorem, spaces of lattices, and geometric approaches to gap problems.

    \item \textbf{Andreas Str\"ombergsson.}  Relevant because of work on
    lattices, dynamics, and the three-gap theorem.

    \item \textbf{Jayadev Athreya.}  Relevant through dynamics, lattices, and
    related gap-distribution questions, although somewhat less directly tied
    to this specific rational cut-and-project problem.

    \item \textbf{Barak Weiss.}  Relevant through homogeneous dynamics,
    Delone sets, and aperiodic order, although the present result may be more
    elementary than his usual focus.
\end{itemize}

\subsection*{Christoffel words, Sturmian sequences, and combinatorics on words}

\begin{itemize}
    \item \textbf{Jean Berstel.}  Central figure in combinatorics on words,
    Sturmian words, and Christoffel words.

    \item \textbf{Christophe Reutenauer.}  Relevant because of major work on
    Christoffel words and combinatorics on words.

    \item \textbf{Aaron Lauve.}  Relevant because of work on Christoffel words
    and related combinatorics.

    \item \textbf{Franco V. Saliola.}  Relevant because of work on Christoffel
    words and combinatorics on words.

    \item \textbf{S\'{e}bastien Labb\'{e}.}  Relevant because of work on
    Christoffel words, multidimensional symbolic dynamics, and cut-and-project
    related structures.

    \item \textbf{Michel Rigo.}  Relevant because of work on combinatorics on
    words, numeration systems, and symbolic dynamics.
\end{itemize}

\section*{Shortlist}

If only a few people are contacted, a sensible order is:

\begin{enumerate}
    \item Alan Haynes.
    \item Henna Koivusalo.
    \item James Walton.
    \item Val\'{e}rie Berth\'{e}.
    \item S\'{e}bastien Labb\'{e}.
    \item Jens Marklof.
\end{enumerate}

The first three are closest to cut-and-project gap questions.  Berth\'{e} and
Labb\'{e} are strong choices if the paper is framed through symbolic dynamics
and combinatorics on words.  Marklof is relevant through the lattice and
gap-problem viewpoint.

\section*{Suggested first-contact email}

\begin{verbatim}
Subject: Short note on rational cut-and-project gap periods

Dear Professor [Name],

I am an independent researcher working on a short note about rational
one-dimensional cut-and-project strip projections.

The paper proves an explicit formula for the minimal period of the projected
multiset gap sequence.  For coprime positive integers alpha, beta and window
parameter omega >= 0, the period is

    lambda = N        if D does not divide N,
    lambda = N / D    if D divides N,

where N = floor(omega alpha) + floor(omega beta) + 1 and
D = alpha^2 + beta^2.

The proof is elementary and reduces the problem to the distribution of
consecutive residue classes modulo D.  The algebraic core has also been
formalised in Lean 4.

Since your work is close to cut-and-project sets / gap problems / symbolic
dynamics, I would be grateful if you were willing to take a brief look at the
manuscript.  In particular, I would appreciate knowing whether the result
appears to be a reasonable arXiv preprint and whether you are aware of prior
literature containing this exact formula.

I fully understand if you do not have time.

Best regards,
Dirk Kunert
\end{verbatim}

\section*{Possible follow-up email asking for endorsement}

This should only be sent if the person has responded positively or has shown
interest in the manuscript.

\begin{verbatim}
Subject: arXiv endorsement question

Dear Professor [Name],

thank you very much for taking the time to look at my manuscript.

I would now like to submit the note to arXiv, probably under [math.CO/math.DS].
Since I am an independent researcher and a first-time submitter in this area,
arXiv requires an endorsement.

If you consider the manuscript suitable as a refereeable mathematical
preprint, would you be willing to endorse me for the corresponding arXiv
category?

Of course I fully understand if you prefer not to do so.

Best regards,
Dirk Kunert
\end{verbatim}

\section*{Checklist before contacting anyone}

Before sending the manuscript to potential readers, check the following:

\begin{enumerate}
    \item The paper compiles without errors.
    \item The title mentions multiplicity or projected multisets.
    \item The abstract clearly states that projected values are counted with
    multiplicity.
    \item The setup contains an explicit multiplicity convention.
    \item The bi-infinite nondecreasing enumeration of projected values is
    justified.
    \item The proof of \(\tilde p_s^{(i+N)}=\tilde p_s^{(i)}+D\) is explicit.
    \item The degenerate case \(D\mid N\) clearly explains the role of zero
    gaps.
    \item The Lean formalisation is described precisely and not overstated.
    \item AI systems are not listed as co-authors.
    \item AI assistance is mentioned only in the acknowledgements.
    \item The bibliography is cleaned up, preferably replacing Wikipedia with
    a mathematical reference for the three-gap theorem.
    \item The GitHub repository builds cleanly if the Lean formalisation is
    advertised.
\end{enumerate}

\section*{Practical recommendation}

The best first step is to contact one or two people from the closest group,
for example Alan Haynes and Henna Koivusalo, with a short email and the
manuscript attached.  If there is no response, proceed to James Walton,
Val\'{e}rie Berth\'{e}, or S\'{e}bastien Labb\'{e}.

The request should initially be for feedback on whether the result appears
suitable for arXiv and whether the formula is known, not immediately for
endorsement.

\end{document}
